% ---------------------------------------------------------------------------
%  Shared preamble: "Electronic Dance Music and Chicago House"
%
%  Engine: XeLaTeX (required -- fontspec + Unicode source text).
%
%  Loaded by every standalone chapter in this directory and by the bound
%  volume edm_guide.tex. Nothing here is chapter-specific; if a rule only
%  applies to one document it belongs in that document, not in this file.
%
%  Font choice follows the Basilisk mini-course build: Cambria for prose,
%  Consolas for the step-grids and file paths. The non-ASCII inventory of the
%  source markdown was audited before this file was written and contains
%  exactly five codepoints (U+00B0, U+00B2, U+00D7, U+00E9, U+2212). Cambria
%  covers all five. The three symbol codepoints are nevertheless mapped to
%  math-font equivalents below, so the build does not silently change
%  appearance if it is ever run on a machine without Cambria installed.
% ---------------------------------------------------------------------------

\usepackage{fontspec}
\setmainfont{Cambria}
\setmonofont{Consolas}[Scale=0.85]

\usepackage{amsmath}
\usepackage{amssymb}
\usepackage{xcolor}
\usepackage{newunicodechar}
\usepackage{graphicx}
\usepackage{array}
\usepackage{enumitem}

% --- symbol codepoints pinned to the math font ----------------------------
% Present in Cambria, mapped anyway so a font substitution cannot quietly
% turn a minus sign into a hyphen in a table of decibel values.
\newunicodechar{−}{\ensuremath{-}}
\newunicodechar{×}{\ensuremath{\times}}
\newunicodechar{²}{\ensuremath{^{2}}}

% --- page geometry, spacing, headers --------------------------------------
\usepackage[letterpaper,margin=0.95in,bottom=1.1in]{geometry}
\usepackage{titlesec}
\usepackage{fancyhdr}

\setlength{\parskip}{0.32em}
\setlength{\parindent}{0pt}
\titlespacing*{\section}{0pt}{1.3em}{0.45em}
\titlespacing*{\subsection}{0pt}{0.9em}{0.3em}

\pagestyle{fancy}
\fancyhf{}
\fancyhead[L]{\small\nouppercase\leftmark}
\fancyhead[R]{\small EDM \& Chicago House}
\fancyfoot[C]{\small\thepage}
\renewcommand{\headrulewidth}{0.4pt}
\fancypagestyle{plain}{\fancyhf{}\fancyfoot[C]{\small\thepage}%
  \renewcommand{\headrulewidth}{0pt}}

% --- verbatim -------------------------------------------------------------
% The step-grids in the rhythm and FL Studio chapters are fixed-width ASCII
% and must not be reflowed or broken. fvextra is loaded for its better
% verbatim handling; breaking is deliberately left off so a 16-step bar
% cannot be silently wrapped into something that is no longer 16 steps.
\usepackage{fvextra}
\fvset{breaklines=false,fontsize=\small}

% --- tables ---------------------------------------------------------------
\usepackage{etoolbox}
\usepackage{booktabs}
\usepackage{longtable}
\AtBeginEnvironment{longtable}{\small}

% --- breaking long identifiers --------------------------------------------
% This guide is full of filenames like 03_first_track.md and sample paths
% like 909 CH 1.wav. TeX will not hyphenate them, so in a narrow table column
% they run into the margin. Permitting a break after each underscore fixes it
% without inserting a hyphen that would misrepresent the filename.
\let\origunderscore\_
\renewcommand{\_}{\origunderscore\allowbreak}

% --- links ----------------------------------------------------------------
\usepackage[colorlinks=true,linkcolor=black!70!blue,urlcolor=black!60!blue,
            citecolor=black!70!blue]{hyperref}

% --- callout boxes --------------------------------------------------------
% Same three-box idiom as the Jamming and Fracture lecture notes, plus a
% fourth for the listening prompts, which are specific to this subject.
\usepackage{mdframed}

\newmdenv[backgroundcolor=blue!5,linecolor=blue!40,linewidth=0.8pt,
  roundcorner=4pt,skipabove=8pt,skipbelow=8pt,
  innertopmargin=6pt,innerbottommargin=6pt,
  innerleftmargin=8pt,innerrightmargin=8pt]{intuitionbox}

\newmdenv[backgroundcolor=orange!8,linecolor=orange!60,linewidth=0.8pt,
  roundcorner=4pt,skipabove=8pt,skipbelow=8pt,
  innertopmargin=6pt,innerbottommargin=6pt,
  innerleftmargin=8pt,innerrightmargin=8pt]{pitfallbox}

\newmdenv[backgroundcolor=gray!10,linecolor=gray!60,linewidth=0.8pt,
  roundcorner=4pt,skipabove=8pt,skipbelow=8pt,
  innertopmargin=6pt,innerbottommargin=6pt,
  innerleftmargin=8pt,innerrightmargin=8pt]{reviewbox}

\newmdenv[backgroundcolor=green!5,linecolor=green!45!black!60,linewidth=0.8pt,
  roundcorner=4pt,skipabove=8pt,skipbelow=8pt,
  innertopmargin=6pt,innerbottommargin=6pt,
  innerleftmargin=8pt,innerrightmargin=8pt]{listenbox}

% --- block quotes ---------------------------------------------------------
\usepackage{quoting}
\renewenvironment{quote}
  {\begin{quoting}[leftmargin=1.2em,vskip=0.4em]\itshape}
  {\end{quoting}}

% --- pandoc compatibility -------------------------------------------------
% The chapter bodies are generated by pandoc, whose LaTeX writer assumes the
% scaffolding its own default template provides. We supply our own template,
% so that scaffolding has to be reproduced here.
%
% \newcounter{none} is the one that bites: pandoc emits every table as an
% unnumbered longtable keyed to a counter of that name, so without it each of
% the ~40 tables in this guide is a hard error, not a warning.
\providecommand{\tightlist}{%
  \setlength{\itemsep}{0pt}\setlength{\parskip}{0pt}}
\providecommand{\passthrough}[1]{#1}

\newcounter{none}
\usepackage{calc}
\makeatletter
\patchcmd\longtable{\par}{\if@noskipsec\mbox{}\fi\par}{}{}
\makeatother
% footnotehyper must come after hyperref; it lets footnotes survive inside a
% longtable head or foot.
\IfFileExists{footnotehyper.sty}{\usepackage{footnotehyper}}{\usepackage{footnote}}
\makesavenoteenv{longtable}

\setlength{\emergencystretch}{3em}
